% automatisch erzeugt von paper/tabellen.py -- nicht von Hand editieren
\begin{table}[htbp]
  \centering
  \footnotesize
  \begin{tabular}{llrrrrrrr}
    \toprule
    Periode & $N$ & \multicolumn{2}{c}{QLIKE} & \multicolumn{2}{c}{MSE ($\times 10^{-7}$)} & \multicolumn{3}{c}{Sharpe} \\
    \cmidrule(lr){3-4}\cmidrule(lr){5-6}\cmidrule(lr){7-9}
     & & Hist. & $\Delta$ DCC & Hist. & DCC & MVP Hist. & MVP DCC & 1/N \\
    \midrule
    2003--2007 & 16--16 & -143,5 & -3,79 & 0,28 & 0,23 & 1,490 & 1,485 & 1,087 \\
    2008--2009 & 16--16 & -119,7 & -11,50 & 21,03 & 16,72 & -0,200 & -0,294 & -0,092 \\
    2010--2019 & 16--17 & -146,1 & -1,75 & 0,79 & 0,60 & 0,940 & 1,036 & 0,760 \\
    2020--2021 & 17--17 & -133,1 & -7,15 & 20,18 & 15,88 & 0,446 & 0,427 & 0,836 \\
    2022--2026 & 17--27 & -205,1 & -0,69 & 1,79 & 1,70 & 0,578 & 0,462 & 0,461 \\
    \bottomrule
  \end{tabular}
  \caption{Teilperioden des Basisfalls: Universum, Prognosequalität und Performance}
  \label{tab:strukturbruch_teilperioden}
  \begin{minipage}{\linewidth}\vspace{0.4em}\centering\footnotesize
  $N$ ist die Spannweite der investierbaren Sektoren; QLIKE-Niveaus sind zwischen Perioden mit unterschiedlichem $N$ nicht vergleichbar.
  \end{minipage}
\end{table}
